\section{Audit of the existing formal proofs}
\label{sec:proof-audit}

The inspected ProveIt revision is
\texttt{b6c7087477e51c017003bd729e712c8969733223}.
The repository already contains three proof-simplification checkpoints.
Consequently, deriving a shared AX-word interpreter, removing the reverse
Lean source certificate, or replacing the original expanded dominated-rule
derivation would repeat work already present in this revision.\cite{verification,phase3}

The current Lean primary dependency closure has 185 project modules,
compared with 204 at the original full-build checkpoint. The latest
phase-three record freshly rebuilds nine affected modules and binds the
176 reused objects to their earlier checked sources, dependencies, objects
and logs. Its two explicit audit clients produce ten reports containing
only \texttt{propext}, \texttt{Classical.choice}, and \texttt{Quot.sound}.
This is an incremental check of the complete dependency closure, not a
fresh compilation of all 185 modules in this session. The original Rocq
record reports 237 selected sources, four closed endpoint assumption
reports and a recursive \texttt{coqchk}. Subsequent simplification records
retain that baseline and identify the affected rebuilt modules.

The precise endpoint in Lean is
\begin{verbatim}
KEndpoint215.blank_halts_iff_zfc_refutation :
  BlankHalts machine215 <-> BProv ZFCax_s [] Form.fBot
\end{verbatim}
The Rocq dense endpoint is
\texttt{tm215\_halts\_iff\_ZFC\_inconsistent} in
\path{k225-formal/rocq/rocq-extension/coq/215/CorrectDense215.v}.
Both refer to the primary machine's literal dense table. The core
composition is already short: the Lean endpoint is a direct application
of the shared \texttt{KEndpoint} theorem, and Rocq rewrites the dense/sparse
equivalence and applies its sparse endpoint. Shortening those final lines
has little mathematical value. The substantial work is in the reusable
simulation and calculus layers, and in the machine-specific certificates.\cite{endpoint}

\begin{center}
\begin{tabular}{lrr}
\hline
Proof source & Original bytes & Current bytes\\
\hline
Lean dominated implication & 689,857 & 35,185\\
Lean AX routine & 202,222 & 7,563\\
Lean source data & 653,569 & 456,548\\
Lean site-coverage facade & 107,243 & 14,916\\
Lean register determinism & 17,209 & 12,988\\
Rocq dominated implication & 921,691 & 22,673\\
\hline
\end{tabular}
\end{center}

These are the repository's recorded source-byte comparisons, rather than
measurements of the size of the resulting kernel proof terms. They also
do not count every retained historical receipt as removed repository data.\cite{simplification}
The current Lean dominated-rule argument has a 61-lemma bootstrap,
81 reused conditional schemas, eight numerical dominance hypotheses, and
a deduction theorem. The 629-node/735-instruction certificate remains
the checked Rocq rendering; it is not the current Lean rendering.

The audit sources are the primary verification README, the simplification
README, the phase-three README and its Lean record, and the actual endpoint
files. Useful current source anchors are
\texttt{KEndpoint215.lean:43} (the endpoint),
\texttt{KSim.lean:317} (positive-step halting transfer),
\texttt{RMOffsetBisim.lean:287} (one certificate gives equivalence),
\texttt{L2AxiomsPerm.lean:100} (AX-word composition), and
\texttt{DomRho8.lean:597} (the unrestricted dominated-row conclusion).
All of these anchors refer to the pinned revision above.

\subsection{A concrete remaining simplification in Rocq}

The Lean source correspondence now uses a single forward certificate and
the mathematical symmetry of an offset bisimulation. The current Rocq
\texttt{Source215.v} still separately evaluates
\begin{verbatim}
certificate_forward : ocheck (ptree p215) (ptree dagPerm) Cert215 = true
certificate_backward :
  ocheck (ptree dagPerm) (ptree p215) (omirror Cert215) = true
\end{verbatim}
with \texttt{vm\_compute} (lines 8--13), and the same directory's
\texttt{OffsetBisim230.v:167}, \texttt{ocheck\_iff}, requires both results.
Thus the same mathematical
improvement remains available in the Rocq lane.

The appropriate addition is the generic theorem
\begin{verbatim}
ocheck_mirror t1 t2 C :
  ocheck t1 t2 C = true ->
  ocheck t2 t1 (omirror C) = true.
\end{verbatim}
Its proof has five small ingredients. Lookup in \texttt{omirror C} is
lookup in \texttt{C} followed by swapping the two paths and negating the
integer offset. The identity
$\delta(-z,i)=-\delta(z,i)$ turns each old count equation into its reverse.
A simple completeness lemma for \texttt{counts\_ok}, paired with the
existing \texttt{counts\_ok\_sound}, avoids reasoning about the order of
the duplicate register indices in its test list. These facts transfer a
branch check, then an entry check, then the initial and universal checks.
The initial offset is zero, so it is fixed by negation. This removes the
second whole-certificate evaluation while preserving the exact theorem
and every table/source input. This Rocq change is a proposed port; the
new files in this packet do not claim to have compiled it.

\subsection{Use a verified parent as the proof boundary}

For a small transition-table edit, the most effective next simplification
is often a direct equivalence proof with the already verified machine.
The entire numerical-calculus calibration, proof enumeration, source
bisimulation, site coverage and blank initializer can then remain behind
one proved interface. The new obligation is a local simulation theorem
about the changed transition graph.

This applies especially well to the new FM215 construction below. The
existing conditional-marker alternative CM215 also has an ordinary
all-tape equivalence with the primary machine, but its invariant needs
an extra fact about state 156 and a two-step incoming path. FM215 uses two
unconditional predecessor-marker facts, each established by the
immediately preceding transition. Its additional correctness proof is
therefore structurally simpler.

This comparison is about the additional proof after the primary endpoint
has been proved. It is not a claim that changing the table removes the
underlying set-theoretic calibration. Likewise, the rebalanced-address
alternative's maximum source depth of 41 can reduce some certificate
arithmetic, but a shorter address is not itself a shorter correctness
argument. The composed allocation shares the primary source and seed,
which makes it a conservative candidate for reusing the existing
certificate pipeline, but it still needs its own physical runtime
certificates unless an additional direct equivalence is established.

\section{A reusable positive simulation theorem}
\label{sec:positive-simulation}

Let $f:X\to X\cup\{\bot\}$ and $g:Y\to Y\cup\{\bot\}$ be deterministic
partial transition functions. A configuration is terminal precisely when
its transition is $\bot$. Let $R\subseteq X\times Y$ satisfy:
\begin{enumerate}
\item If $R(x,y)$ and $f(x)=x'$, then some $y'$ satisfies $R(x',y')$
  and $y\mathrel{\longrightarrow_g^{k}}y'$ for some $k\geq1$.
\item If $R(x,y)$ and $f(x)=\bot$, then the $g$-computation from $y$
  halts after finitely many steps.
\end{enumerate}
Then $R(x,y)$ implies that $f$ halts from $x$ if and only if $g$ halts
from $y$. The relation need not be injective, surjective, or functional.
Reachability invariants can simply be included in $R$.

For the forward implication, induct on the finite source computation,
concatenating the finitely many target blocks. For the reverse implication,
induct on the number $t$ of successful target steps before termination.
If the source is terminal, the conclusion is immediate. Otherwise its
next step is simulated by a block of length $k\geq1$. This block cannot
extend past the target's terminal configuration: determinism and the
absence of a transition there would contradict its certified successful
completion. Hence $k\leq t$, and the induction hypothesis applies at the
next related boundary with the strictly smaller bound $t-k$.

Positive progress is essential. If zero-length target blocks were
allowed, a nonhalting source self-loop could be related to a terminal
target forever. No uniform upper bound on block lengths is needed for
halting equivalence. In FM215 there is nevertheless the useful exact
bound $k\in\{1,3\}$.

The supplied \texttt{PositiveSimulation.lean} is a generic extraction of
the final argument in the repository's \texttt{KSim.halts\_iff}. It uses
arbitrary configuration types and deterministic partial transitions, so
it can also serve machine-to-machine optimizations. It is not presented
as a new abstract computability theorem. Its value is as a small reusable
formal interface. The accompanying execution receipt records a successful
Lean 4.32.0 compilation of this theorem. Its printed dependencies are
exactly \texttt{propext} and \texttt{Quot.sound}.

\section{FM215: the short correctness argument}
\label{sec:fm-proof}

Write $T$ for the primary native K215 table and $N$ for FM215. Define the
state map
\[
 \phi(127)=124,\qquad \phi(q)=q\quad(q\neq127).
\]
The two relevant old states have rows
\[
\begin{array}{c|cc}
 &0&1\\\hline
124&0R62&1R138\\
127&0R14&1R138.
\end{array}
\]
Their entire incoming-edge census is
$104/1=1R124$ and $110/0=0R127$. Hence, on every old computation from
initial control zero and on every initial tape,
\[
 q=124\Longrightarrow u(h-1)=1,\qquad
 q=127\Longrightarrow u(h-1)=0,
\]
where $u$ denotes the current tape and $h$ the current head position.
Call this two-clause assertion $\mathcal I$. It holds initially because control
zero is neither exceptional state, and it holds after \emph{every}
old transition: the two possible destination states force their marker
through their unique incoming edge. In particular its preservation
requires no hypothesis about the preceding tape.

FM215 changes exactly four rows:
\[
\begin{array}{c|cc}
\text{row}&\text{primary K215}&\text{FM215}\\\hline
124/0&0R62&0L127\\
127/0&0R14&0R8\\
127/1&1R138&1R109\\
110/0&0R127&0R124.
\end{array}
\]
The existing provider rows are $8/0=0R14$ and $109/0=0R62$.
At simulation boundaries relate configurations with the same entire tape
and head, old invariant $\mathcal I$, and new control $\phi(q)$.
An ordinary old transition is simulated by one new transition. On the
common read-one branch, both old exceptional states go to 138, and their
new representative 124 does the same.

On an exceptional read-zero branch, let the predecessor marker be $m$.
The new computation is
\[
 (124,h)\longrightarrow(127,h-1)
 \longrightarrow(x_m,h)\longrightarrow(a_m,h+1),
\]
where $(x_0,a_0)=(8,14)$ and $(x_1,a_1)=(109,62)$.
The three writes are respectively $0,m,0$, each equal to the value
already present. Thus the entire tape is unchanged. The invariant picks
the correct old destination, and the head has moved right exactly once.
All three internal controls have a transition, so no halt is hidden in
the block. The state map fixes the terminal label 254.

The positive simulation theorem now proves halting equivalence on all
binary initial tapes, with any initial head position, from control zero.
If $z(t)$ counts the old read-zero visits to states 124 or 127 in a
successful prefix of length $t$, the exact corresponding new duration is
\[
 t+2z(t),\qquad t\leq t+2z(t)\leq3t.
\]
This is a statement about arbitrary whole tapes, not a finite sample of
register values and not an inference from the blank prefix.

\subsection{How the formal certificate binds the machine}

\texttt{FM215Data.lean} embeds both complete native tables and defines
functional tapes $\mathbb Z\to\{0,1\}$. The ambient label type is
\texttt{Fin 512}, because the native tables use labels up to 302. This
is an indexing device: exactly 215 labels have working rows, and every
outgoing target is a working label or the separate terminal label 254.
The finite ambient type does not assert a 512-state machine.

\texttt{FM215.lean}'s \texttt{incoming} theorem checks the two required
incoming-row properties for every ambient state and both symbols.
\texttt{outer} checks all unchanged and redirected rows under $\phi$.
Both are finite kernel decisions. The subsequent proofs establish
invariant preservation, the complete three-step tape identity, and the
one-or-three-step simulation. The endpoint is
\begin{verbatim}
FM215.all_tapes_halts_iff (h : Int) (t : Tape) :
  Halts (step old) (0,h,t) <-> Halts (step new) (0,h,t).
\end{verbatim}
The actual file uses the corresponding structure notation and qualified
\path{PositiveSimulation.Halts}. The blank instance is obtained by
$h=0$ and the constantly false tape.

All three supplied modules were compiled successfully in this session
with the official Lean 4.32.0 release, commit
\texttt{8c9756b28d64dab099da31a4c09229a9e6a2ef35}. The generic theorem
and both FM215 endpoints report exactly \texttt{propext} and
\texttt{Quot.sound}; there is no admitted theorem, extra assumption, or
native evaluation oracle. The final compiler logs contain no warnings.
\texttt{fm215-kernel-receipt.json} binds the checked sources, compiled
objects and logs by SHA-256. The successful sequential run took about
26 seconds in this environment; this is a reproducibility datum rather
than a comparative benchmark. The generic proof file has 99 lines,
the FM215 proof file 172, and the table-and-semantics file 918, counting
comments and blank lines. The first two files together are 10,375 bytes.
\texttt{check\_fm215.py} reproduces the compilation and the binding audit
without any dependency on the ProveIt project build.

The independent \texttt{check\_lean\_bindings.py} reads every literal Lean
row and compares it with the exact \texttt{.tm} input. It also checks the
ambient bound, so a large numerical literal cannot silently wrap under
\texttt{Fin}'s numeral instance. The primary native SHA-256 is
\begin{verbatim}
9683d6ed9628e6ce407c315b08f5da3d063da94e92706f275ca0f33e8272e4b2
\end{verbatim}
and the FM215 native SHA-256 is
\begin{verbatim}
3ece1233f11d33375425221c5bcedba941e073763912990eabf321197772dce5
\end{verbatim}
The data audit is separate from the Lean compiler check. An exact-table
theorem in functional-tape semantics is also separate from composing the
repository's existing \texttt{KEndpoint215} theorem, whose interface uses
its own dense and zipper semantics. The outstanding imported bridge has
two specific parts: identify the old native table with
\texttt{KTable215.lean}'s \texttt{tm215} on working and terminal controls
under the repository's renaming $254\mapsto230$, $301\mapsto102$,
$302\mapsto120$; and relate functional tapes to \texttt{KZip}'s two
finite lists, proving the one-step and terminal cases. The existing
\texttt{KDenseBridge} already supplies the dense-to-zipper side.
Alternatively one can adapt that generic bridge directly from Core's
finite-support tapes to functional tapes, using the native dense
relabelling. Then the new blank equivalence composes with
\texttt{KEndpoint215}. An endpoint for the new dense table also uses
FM215's supplied canonicalization map. These bridges must be explicit
if a new end-to-end Lean consistency endpoint is claimed. No new blank
bootstrap or calculus derivation is needed. Ordinary mathematical
composition with the established primary calibration already gives
the desired ZFC equivalence.

\subsection{Independent check of the existing CM215 route}

For comparison, \texttt{check\_cm215.py} in this packet independently checks
the existing six-row conditional-marker alternative against the exact
primary bytes. Its invariant is
\[
 q=52\Rightarrow u(h-1)=1,\quad
 (q=176\ \&\ u(h)=0)\Rightarrow u(h-1)=0,\quad
 q=156\Rightarrow u(h+1)=1.
\]
An old radius-two tape window determines the radius-one invariant after
one L/R transition. Enumerating all 32 such windows per working control
is therefore an exact finite proof-obligation check for arbitrary tapes.
Of 6,880 state/window combinations, 6,840 satisfy the invariant. All
preserve it and pass the whole-window block identity; 6,824 use one step
and 16 use three. The checker also verifies the supplied native/dense
bijection. It does not import the earlier symbolic runtime checker,
the earlier blank trace, or the source calculus. The ordinary locality
argument lifts these finite obligations to whole tapes because every
read and write lies inside the certified window and exterior cells are
untouched. The simpler FM215 invariant is the reason to prefer its new
formal proof as the next compact endpoint.
